\begin{mdframed}
    \textbf{La extensión máxima para esta sección es de 4 páginas.} Si utiliza menos páginas y explica bien sus resultados, es preferible.
\end{mdframed}

\subsection{Sorting}

La Figura~\ref{fig:sorting_tiempo} muestra el tiempo de ordenamiento en función del tamaño del arreglo. Se observa un crecimiento consistente del tiempo al aumentar $n$. En particular, \textit{quick sort} presenta el mejor desempeño general, seguido por \texttt{std::sort}, mientras que \textit{merge sort} presenta los tiempos más altos.

\begin{figure}[H]
    \centering
    \includegraphics[width=0.8\textwidth]{tiempo_ms_vs_n.png}
    \caption{Tiempo de ordenamiento vs tamaño del arreglo.}
    \label{fig:sorting_tiempo}
\end{figure}

La Figura~\ref{fig:sorting_loglog} muestra los mismos resultados en escala logarítmica. Se observa una tendencia aproximadamente lineal, consistente con el comportamiento teórico esperado.

\begin{figure}[H]
    \centering
    \includegraphics[width=0.8\textwidth]{tiempo_ms_loglog.png}
    \caption{Tiempo de ordenamiento en escala log-log.}
    \label{fig:sorting_loglog}
\end{figure}

La Figura~\ref{fig:sorting_memoria} muestra el uso de memoria. Se observa que \textit{merge sort} presenta un mayor crecimiento en memoria en comparación con los otros algoritmos.

\begin{figure}[H]
    \centering
    \includegraphics[width=0.8\textwidth]{memoria_kb_vs_n.png}
    \caption{Uso de memoria vs tamaño del arreglo.}
    \label{fig:sorting_memoria}
\end{figure}


\subsection{Multiplicación de matrices}

La Figura~\ref{fig:matrices_tiempo} muestra el tiempo de ejecución en función del tamaño. Ambas implementaciones aumentan su tiempo al crecer $n$, aunque la diferencia entre ellas no es extremadamente marcada.

\begin{figure}[H]
    \centering
    \includegraphics[width=0.8\textwidth]{tiempo_ms_vs_n_matrix.png}
    \caption{Tiempo de multiplicación de matrices vs tamaño.}
    \label{fig:matrices_tiempo}
\end{figure}

La Figura~\ref{fig:matrices_loglog} muestra los resultados en escala logarítmica. Ambas curvas presentan un comportamiento similar.

\begin{figure}[H]
    \centering
    \includegraphics[width=0.8\textwidth]{tiempo_ms_loglog_matrix.png}
    \caption{Tiempo en escala log-log para multiplicación de matrices.}
    \label{fig:matrices_loglog}
\end{figure}

La Figura~\ref{fig:matrices_memoria} muestra el uso de memoria. Strassen utiliza más memoria que la versión naive.

\begin{figure}[H]
    \centering
    \includegraphics[width=0.8\textwidth]{memoria_kb_vs_n_matrix.png}
    \caption{Uso de memoria en multiplicación de matrices.}
    \label{fig:matrices_memoria}
\end{figure}
